% Unit 126 terjemahan Bahasa Indonesia dari chapter9.tex baris 1043--1156.
% Sumber: Methods of Algebra, Volume 2, commit
% 9a5803ff77dd3257484cb177f851a73770a59dd3 (CC BY 4.0).
% stable-unit-id: o014.aljabr2.chapter9.coendomorphism-coalgebra-a
% Cakupan: Bagian 9.4 dari sifat universal coHom hingga konvensi kurung.
% Koreksi sumber terpaksa: baris 1105 membalik omega_1 dan omega_2 pada
% koaksi kanonik; terjemahan memakai omega_2(TX)->omega_1(TX) tensor coHom.

% segment-id: o014.li2.chapter9.coendomorphism-coalgebra-a.h001
\section{Koaljabar endomorfisme}\label{sec:coend}
\hypertarget{o014.aljabr2.chapter9.coendomorphism-coalgebra-a}{ }
% segment-id: o014.li2.chapter9.coendomorphism-coalgebra-a.p001
Mulai bagian ini, kita menyiapkan serangkaian landasan bagi teori kategori
Tannaka (lihat \S\ref{sec:Tannakian-cat}); beberapa alat dan gagasan yang
terlibat memiliki penerapan yang lebih luas. Rujukan utama ialah
\cite{Del90}.

% segment-id: o014.li2.chapter9.coendomorphism-coalgebra-a.p002
Misalkan $\Bbbk$ gelanggang komutatif. Untuk kategori $\mathcal{A}$ dan
funktor $\xi:\mathcal{A}\to\Bbbk\dcate{Mod}$ sembarang, kita dapat meninjau
aljabar-$\Bbbk$ endomorfisme funktor, $\End(\xi)$. Aljabar ini memberikan
keluarga pemetaan linear-$\Bbbk$
% segment-id: o014.li2.chapter9.coendomorphism-coalgebra-a.q001
\[ \End(\xi) \dotimes{\Bbbk} \xi(X) \to \xi(X), \quad X \in \Obj(\mathcal{A}), \]
yang funktorial dalam $X$, atau secara ringkas
$\End(\xi)\dotimes{\Bbbk}\xi\to\xi$. Selain itu, $\End(\xi)$ universal
untuk tindakan semacam ini: secara tautologis, memberikan tindakan suatu
aljabar-$\Bbbk$ $E$ pada $\xi$ sama artinya dengan memberikan
homomorfisme aljabar $E\to\End(\xi)$. Struktur perkalian $\End(\xi)$ dapat
diturunkan secara abstrak dari sifat universal ini.

% segment-id: o014.li2.chapter9.coendomorphism-coalgebra-a.p003
Untuk persoalan yang akan kita tangani, versi dual konstruksi ini lebih
berguna sekaligus lebih sederhana. Namun kita perlu meninjau funktor yang
mengambil nilai dalam $\cated{Mod}B$, dengan $B$ sebuah aljabar-$\Bbbk$
yang, menurut konvensi, tak nol. Konstruksi tersebut berkaitan erat dengan
konstruksi abstrak yang disebut \emph{koend} dalam teori kategori; latihan
bab ini memberikan penjelasan lebih lanjut.

% segment-id: o014.li2.chapter9.coendomorphism-coalgebra-a.d001
\begin{definition}\label{def:cogebra-L-universal}
	\index[sym1]{coHom-coEnd@$\coHom, \coEnd$}
	Misalkan $\Bbbk$ gelanggang komutatif, $B$ aljabar-$\Bbbk$, dan
	$\omega_1,\omega_2:\mathcal{A}\to\cated{Mod}B$ funktor. Tinjau suatu
	bimodul-$(B,B)$ $L$ beserta transformasi alami
	$a:\omega_2\to\omega_1\dotimes{B}L$, yakni keluarga homomorfisme
	modul-$B$ kanan
	% segment-id: o014.li2.chapter9.coendomorphism-coalgebra-a.q002
	\[ a_X: \omega_2(X) \to \omega_1(X) \dotimes{B} L, \quad X \in \Obj(\mathcal{A}), \]
	sedemikian rupa sehingga diagram berikut komutatif untuk setiap
	$f\in\Hom_{\mathcal{A}}(X,Y)$:
	% segment-id: o014.li2.chapter9.coendomorphism-coalgebra-a.q003
	\[\begin{tikzcd}
		\omega_2(X) \arrow[r, "a_X"] \arrow[d, "{\omega_2(f)}"'] & \omega_1(X) \dotimes{B} L \arrow[d, "{\omega_1(f) \otimes \identity}"] \\
		\omega_2(Y) \arrow[r, "a_Y"'] & \omega_1(Y) \dotimes{B} L.
	\end{tikzcd}\]

	Sebuah morfisme dari data $(L,a)$ ke $(L',a')$ didefinisikan sebagai
	homomorfisme bimodul $t:L\to L'$ yang membuat diagram berikut selalu
	komutatif:
	% segment-id: o014.li2.chapter9.coendomorphism-coalgebra-a.q004
	\[\begin{tikzcd}
		\omega_2(X) \arrow[r, "{a_X}"] \arrow[rd, "{a'_X}"'] & \omega_1(X) \dotimes{B} L \arrow[d, "{\identity \otimes t}"] \\
		& \omega_1(X) \dotimes{B} L'.
	\end{tikzcd}\]
	Sebaliknya, data $(L,a)$ dan homomorfisme bimodul $t:L\to L'$
	menentukan $a'$ melalui diagram ini, beserta morfisme
	$t:(L,a)\to(L',a')$.

	Jika kategori semua data $(L,a)$ dan morfisme tersebut memiliki objek
	awal, bimodul-$(B,B)$ yang bersesuaian dinotasikan oleh
	% segment-id: o014.li2.chapter9.coendomorphism-coalgebra-a.q005
	\[ \coHom(\omega_1, \omega_2) = \coHom_{\Bbbk}(\omega_1, \omega_2), \]
	dan morfisme yang bersesuaian dinotasikan
	$\lambda:\omega_2\to\omega_1\dotimes{B}L$. Kita juga menuliskan
	% segment-id: o014.li2.chapter9.coendomorphism-coalgebra-a.q006
	\[ \coEnd(\omega) = \coEnd_{\Bbbk}(\omega) := \coHom(\omega, \omega). \]
\end{definition}

% segment-id: o014.li2.chapter9.coendomorphism-coalgebra-a.p004
Simbol $\coHom$ berarti $\operatorname{coHom}$ dan $\coEnd$ berarti
$\operatorname{coEnd}$; keduanya harus dipandang sebagai dual $\Hom$ dan
$\End$, dan singkatannya hanya untuk tata letak. Perhatikan bahwa konsep
bimodul-$(B,B)$ tidak hanya bergantung pada struktur gelanggang $B$, tetapi
juga pada $\Bbbk$.

% segment-id: o014.li2.chapter9.coendomorphism-coalgebra-a.r001
\begin{remark}
	Persamaan $\coHom(\omega_1,\omega_2)=0$ berarti bahwa setiap morfisme
	$\omega_2\to\omega_1\dotimes{B}L$ adalah nol; definisi tidak menutup
	kemungkinan ini. Sebaliknya, jika $\coEnd(\omega)$ ada dan $\omega$
	bukan funktor nol konstan, maka $\coEnd(\omega)\neq0$, sebab
	$\omega\rightiso\omega\dotimes{B}B$.
\end{remark}

% segment-id: o014.li2.chapter9.coendomorphism-coalgebra-a.p005
Jika $\coEnd(\omega)$ ada, menerapkan sifat universal pada
$\omega\rightiso\omega\dotimes{B}B$ memberikan homomorfisme bimodul
% segment-id: o014.li2.chapter9.coendomorphism-coalgebra-a.q007
\[ \epsilon: \coEnd(\omega) \to B. \]

% segment-id: o014.li2.chapter9.coendomorphism-coalgebra-a.p006
Di sisi lain, jika semua $\coHom$ yang dibahas ada, penerapan berurutan
pada $\omega_1,\omega_2,\omega_3$ menghasilkan morfisme
% segment-id: o014.li2.chapter9.coendomorphism-coalgebra-a.q008
\begin{equation*}
	\omega_3 \to \omega_2 \dotimes{B} \coHom(\omega_2, \omega_3) \to \omega_1 \dotimes{B} \coHom(\omega_1, \omega_2) \dotimes{B} \coHom(\omega_2, \omega_3).
\end{equation*}
Sifat universal selanjutnya memberikan homomorfisme bimodul kanonik
% segment-id: o014.li2.chapter9.coendomorphism-coalgebra-a.q009
\begin{equation}\label{eqn:L-comult-prep}
	\coHom(\omega_1, \omega_3) \to \coHom(\omega_1, \omega_2) \dotimes{B} \coHom(\omega_2, \omega_3).
\end{equation}

% segment-id: o014.li2.chapter9.coendomorphism-coalgebra-a.p007
Pemeriksaan langsung dengan funktor $\omega_1,\ldots,\omega_4$ menunjukkan bahwa
\eqref{eqn:L-comult-prep} memenuhi koasosiativitas. Karena itu
$\coEnd(\omega)$ menjadi koaljabar dalam $(B,B)\dcate{Mod}$, dengan
$\epsilon:\coEnd(\omega)\to B$ sebagai kounit, sedangkan setiap
$\omega(X)$ menjadi komodul-$\coEnd(\omega)$ melalui morfisme kanonik
$\lambda$. Pembahasan ini dirangkum sebagai berikut.

% segment-id: o014.li2.chapter9.coendomorphism-coalgebra-a.dp001
\begin{definition-proposition}[Koaljabar endomorfisme]\label{def:cogebra-L}
	\index{koaljabar endomorfisme@koaljabar endomorfisme (endomorphism coalgebra)}
	Ambil funktor $\omega:\mathcal{A}\to\cated{Mod}B$ dan andaikan
	$\coEnd(\omega)$ ada.
	\begin{itemize}
		\item Sebagai objek kategori monoidal $(B,B)\dcate{Mod}$,
		$\coEnd(\omega)$ memiliki struktur koaljabar kanonik sedemikian rupa
		sehingga setiap $\omega(X)$ menjadi komodul-$\coEnd(\omega)$ melalui
		morfisme kanonik, dan morfisme dalam $\mathcal{A}$ menginduksi
		homomorfisme komodul. Kita menyebut $\coEnd(\omega)$
		\emph{koaljabar endomorfisme} dari $\omega$.\footnote{Mungkin nama
		yang lebih tepat ialah koaljabar koendomorfisme.}
		\item Untuk setiap koaljabar $L$ dalam $(B,B)\dcate{Mod}$,
		memberikan struktur komodul-$L$ pada semua $\omega(X)$ secara
		kompatibel sama artinya dengan memberikan homomorfisme koaljabar
		$\coEnd(\omega)\to L$.
	\end{itemize}
\end{definition-proposition}

% segment-id: o014.li2.chapter9.coendomorphism-coalgebra-a.r002
\begin{remark}\label{rem:L-cogebra-functoriality}
	Untuk suatu funktor $T:\mathcal{A}'\to\mathcal{A}$ dan
	$\omega_1,\omega_2:\mathcal{A}\to\cated{Mod}B$, jika $\coHom$ yang
	dibahas ada, morfisme kanonik
	$\omega_2(TX)\to\omega_1(TX)\dotimes{B}\coHom(\omega_1,\omega_2)$
	bersama sifat universal secara alami menginduksi homomorfisme
	% segment-id: o014.li2.chapter9.coendomorphism-coalgebra-a.q010
	\[ \coHom(\omega_1 T, \omega_2 T) \to \coHom(\omega_1, \omega_2), \]
	yang kompatibel dengan semua struktur di atas. Khususnya, untuk
	$\omega_1=\omega=\omega_2$ diperoleh homomorfisme koaljabar
	$\coEnd(\omega T)\to\coEnd(\omega)$.
\end{remark}

% segment-id: o014.li2.chapter9.coendomorphism-coalgebra-a.pr001
\begin{proposition}\label{prop:coHom-colim}
	Misalkan
	$\omega_1,\omega_2:\mathcal{A}\to\cated{Mod}B$ funktor. Andaikan
	$\mathcal{A}$ merupakan gabungan suatu keluarga subkategori, yang masing-masing
	ditulis $\mathcal{A}'$, dan tuliskan
	$\omega_i':=\omega_i|_{\mathcal{A}'}$. Jika
	$\coHom(\omega_1',\omega_2')$ ada untuk setiap $\mathcal{A}'$, maka
	$\coHom(\omega_1,\omega_2)$ ada dan terdapat isomorfisme kanonik
	% segment-id: o014.li2.chapter9.coendomorphism-coalgebra-a.q011
	\[ \varinjlim_{\mathcal{A}'} \coHom(\omega'_1, \omega'_2) \rightiso \coHom(\omega_1, \omega_2). \]
	Kolimit di ruas kiri diurutkan oleh inklusi; homomorfisme transisinya
	diberikan seperti dalam Catatan~\ref{rem:L-cogebra-functoriality}.
	Isomorfisme tersebut kompatibel dengan \eqref{eqn:L-comult-prep}.
\end{proposition}
\begin{proof}
	Untuk setiap bimodul-$(B,B)$ $L$, memberikan morfisme
	$\omega_2\to\omega_1\dotimes{B}L$ sama artinya dengan memberikan
	$\omega_2'\to\omega_1'\dotimes{B}L$ untuk setiap $\mathcal{A}'$ secara
	kompatibel; ini juga sama artinya dengan memberikan keluarga kompatibel
	homomorfisme $\coHom(\omega_1',\omega_2')\to L$.
\end{proof}

% segment-id: o014.li2.chapter9.coendomorphism-coalgebra-a.p008
Persoalannya ialah menjamin keberadaan $\coHom(\omega_1,\omega_2)$, yang
tidak serta-merta. Mula-mula kita harus mengendalikan ukuran
kategori. Pilih semesta Grothendieck $\mathcal{U}$, sehingga konsep
himpunan kecil dan kategori kecil terdefinisi. Sebuah kategori disebut
\emph{kecil secara esensial} jika kelas isomorfisme objek-objeknya
membentuk himpunan kecil (Definisi~\ref{def:essentially-small-cat}).

% segment-id: o014.li2.chapter9.coendomorphism-coalgebra-a.cv001
\begin{convention}\label{con:Mod-pfg}
	\index[sym1]{Modpfg@$\cate{Mod}_{\mathrm{pfg}}$}
	Misalkan $\Bbbk$ gelanggang komutatif dan $B$ aljabar-$\Bbbk$. Kita
	memerlukan notasi berikut.
	\begin{itemize}
		\item Tinjau subkategori penuh
		% segment-id: o014.li2.chapter9.coendomorphism-coalgebra-a.q012
		\[ \underbracket{\catesubd{Mod}{\mathrm{pfg}}B}_{\text{projektif + terbangkit hingga}} \subset \underbracket{\catesubd{Mod}{\mathrm{fg}}B}_{\text{terbangkit hingga}} \subset \cated{Mod}B. \]
		\item Untuk setiap modul-$B$ kanan $P$, definisikan modul-$B$ kiri
		$P^\vee$ menurut \eqref{eqn:P-vee}. Setiap homomorfisme
		$\varphi:P\to Q$ menginduksi homomorfisme dual
		$\varphi^\vee:Q^\vee\to P^\vee$.
	\end{itemize}
	Menurut konvensi buku ini tentang struktur aljabar, $\Bbbk$ dan $B$
	direalisasikan pada himpunan kecil. Syarat terbangkit hingga memastikan
	bahwa $\catesubd{Mod}{\mathrm{fg}}B$ dan
	$\cate{Vect}_{\mathrm f}(\Bbbk)$ kecil secara esensial; keduanya juga
	$\Bbbk$-linear. Namun kategori kecil secara esensial $\mathcal{A}$ yang
	dibahas dalam bagian ini tidak harus memiliki struktur $\Bbbk$-linear.
\end{convention}

% segment-id: o014.li2.chapter9.coendomorphism-coalgebra-a.p009
Untuk setiap
$P_1,P_2\in\Obj(\catesubd{Mod}{\mathrm{pfg}}B)$, terdapat isomorfisme
kanonik
% segment-id: o014.li2.chapter9.coendomorphism-coalgebra-a.q013
\begin{equation*}\begin{aligned}
	\Hom_{\cated{Mod}B}\left( P_2, P_1 \dotimes{B} L \right) & \rightiso \Hom_{(B, B)\dcate{Mod}}\left( P_1^\vee \dotimes{\Bbbk} P_2, L \right) \\
	\varphi & \mapsto \left[ \check{p}_1 \otimes p_2 \mapsto (\check{p}_1 \otimes \identity_L)(\varphi(p_2)) \right] .
\end{aligned}\end{equation*}
Hal ini mudah direduksi ke kasus $P_1$ dan $P_2$ modul bebas berperingkat
hingga. Jadi, bila
$\omega_1,\omega_2:\mathcal{A}\to\cated{Mod}B$ mengambil nilai dalam
$\catesubd{Mod}{\mathrm{pfg}}B$, memberikan data $a=(a_X)_X$ dari
Definisi~\ref{def:cogebra-L-universal} setara dengan memberikan keluarga
homomorfisme
% segment-id: o014.li2.chapter9.coendomorphism-coalgebra-a.q014
\[ \alpha_X: \omega_1(X)^\vee \dotimes{\Bbbk} \omega_2(X) \to L, \quad X \in \Obj(\mathcal{A}), \]
dengan funktorialitas dalam $X$ yang diterjemahkan sebagai komutativitas
diagram berikut untuk setiap $f\in\Hom_{\mathcal{A}}(X,Y)$:
% segment-id: o014.li2.chapter9.coendomorphism-coalgebra-a.q015
\begin{equation}\label{eqn:cogebra-L}\begin{tikzcd}
	\omega_1(Y)^\vee \dotimes{\Bbbk} \omega_2(X) \arrow[r, "{\omega_1(f)^\vee \otimes \identity}" inner sep=0.6em] \arrow[d, "{\identity \otimes \omega_2(f)}"'] & \omega_1(X)^\vee \dotimes{\Bbbk} \omega_2(X) \arrow[d, "{\alpha_X}"] \\
	\omega_1(Y)^\vee \dotimes{\Bbbk} \omega_2(Y) \arrow[r, "{\alpha_Y}"'] & L.
\end{tikzcd}\end{equation}

% segment-id: o014.li2.chapter9.coendomorphism-coalgebra-a.pr002
\begin{proposition}\label{prop:cogebra-L-universal}
	Misalkan $\mathcal{A}$ kategori kecil secara esensial dan
	$\omega_1,\omega_2:\mathcal{A}\to
	\catesubd{Mod}{\mathrm{pfg}}B$ funktor sembarang. Maka
	$\coHom(\omega_1,\omega_2)$ dari
	Definisi~\ref{def:cogebra-L-universal} ada.
\end{proposition}
\begin{proof}
	Satu-satunya hal yang perlu dibuktikan ialah keberadaan objek awal.
	Gunakan diagram komutatif \eqref{eqn:cogebra-L}. Ambil $L_0$ sebagai
	jumlah langsung semua
	$\omega_1(X)^\vee\dotimes{\Bbbk}\omega_2(X)$, dengan $X$ berkisar atas
	suatu keluarga wakil $\Obj(\mathcal{A})/\simeq$; di sinilah syarat kecil
	secara esensial diperlukan. Untuk setiap $f:X\to Y$, definisikan
	% segment-id: o014.li2.chapter9.coendomorphism-coalgebra-a.q016
	\[ \delta_f := \omega_1(f)^\vee \otimes \identity - \identity \otimes \omega_2(f) : \omega_1(Y)^\vee \dotimes{\Bbbk} \omega_2(X) \to L_0. \]
	Kemudian
	$\coHom(\omega_1,\omega_2):=L_0/\sum_f\Image(\delta_f)$ adalah objek
	yang dicari.
\end{proof}

% segment-id: o014.li2.chapter9.coendomorphism-coalgebra-a.cv002
\begin{convention}\label{con:coend-bracket}
	Notasi konkret berikut akan berguna. Untuk
	$t\in\omega_1(X)^\vee\dotimes{\Bbbk}\omega_2(X)$, tuliskan $[t]$ untuk
	bayangannya dalam $\coHom(\omega_1,\omega_2)$; bayangan-bayangan ini
	membangkitkan $\coHom(\omega_1,\omega_2)$. Dalam sifat universal
	Definisi~\ref{def:cogebra-L}, bayangan
	$[\phi\otimes x]\in\coEnd(\omega)$ dalam $L$ diperoleh dengan
	memetakan $x\in\omega(X)$ melalui
	$\omega(X)\to\omega(X)\dotimes{B}L$, lalu mengontraksikannya dengan
	$\phi\in\omega(X)^\vee$.
\end{convention}
